% Section 4.4: Breaking the Stalemate
% Transition section: motivates the globular cluster analysis (Paper 3)

\section{Breaking the Stalemate: Independent Constraints from Globular Clusters}
\label{sec:4.4}

% The analyses discussed in the preceding sections share a common vulnerability: they all operate directly within the complex gamma-ray environment of the Galactic Center, where the results depend sensitively on the adopted interstellar emission model, masking procedure, and treatment of the Fermi Bubbles.
\blue{The analyses discussed in the preceding sections share a common vulnerability: they all operate directly within the complex gamma-ray environment of the Galactic Center, where the results depend sensitively on the diffuse-emission systematics discussed in Section~\ref{sec:4.3}.}
% Whether one uses template fitting, non-Poissonian statistics, or neural network-based inference, the systematic uncertainties associated with the Galactic diffuse emission propagate into every conclusion about the nature of the GCE.
\blue{Whether one uses template fitting, non-Poissonian statistics, or neural network-based inference, every conclusion about the nature of the GCE inherits these uncertainties.}

To make further progress, independent constraints from environments that do not suffer from these systematics are needed.

Globular clusters are well suited for this purpose.
% The high stellar densities in these systems drive frequent dynamical interactions, producing large populations of low-mass X-ray binaries and, through recycling, millisecond pulsars~\cite{Amerio:2024qor}.
% The number of MSPs in a given cluster correlates with observable macroscopic properties --- the stellar encounter rate $\Gamma_e$ and the visible-band luminosity $L_V$ (a proxy for stellar mass) --- allowing the total MSP content to be estimated from quantities that are measured independently of the gamma-ray data~\cite{Amerio:2024qor}.
\blue{Globular clusters host large, dynamically formed populations of millisecond pulsars, and the MSP content of a given cluster correlates with macroscopic properties --- the stellar encounter rate and the visible-band luminosity --- that are measured independently of the gamma-ray data (Section~\ref{sec:msp_paper})~\cite{Amerio:2024qor}.}
Because globular clusters are discrete, spatially isolated systems far from the confused Galactic plane, their gamma-ray emission can be measured cleanly, without the severe interstellar emission model uncertainties that affect analyses of the GCE.

The age of the stellar populations in globular clusters is especially relevant to the GCE debate.
% The MSP hypothesis requires the pulsars responsible for the excess to be systematically fainter than those observed in the Galactic Plane, since otherwise Fermi-LAT should have already resolved many of them individually~\cite{Hooper:2016rap,Holst:2024fvb}.
\blue{The MSP hypothesis requires the pulsars responsible for the excess to be systematically fainter than those observed in the Galactic Plane, since otherwise Fermi-LAT should have already resolved many of them individually (cf.\ Section~\ref{sec:4.2.2}).}
One natural explanation is that the Galactic Bulge formed its stars early, so its pulsars are very old and have lost much of their rotational kinetic energy, becoming intrinsically dimmer~\cite{Amerio:2024qor}.
Globular clusters test this hypothesis directly: their stellar populations are comparably old to the Bulge, and if extreme age naturally produces faint MSPs, that faintness should be observable in these systems as well.

% Building on an earlier study by Hooper and Linden~\cite{Hooper:2016rap} that used 7 years of data, in Amerio, Hooper, and Linden~\cite{Amerio:2024qor} we measured the gamma-ray luminosity function of MSPs across the Milky Way's globular cluster system using 15.8 years of Fermi-LAT data.
% We detected statistically significant gamma-ray emission from 56 of the 157 known globular clusters, 8 of which had not appeared in previous source catalogs.
% By fitting the observed gamma-ray luminosities of 87 clusters against predictions based on their stellar encounter rates and visible luminosities, we constrained the underlying log-normal luminosity function.
% The analysis consistently favors a mean pulsar luminosity of $\langle L_\gamma \rangle \sim (1{-}8) \times 10^{33}$~erg/s (integrated between 0.1 and 100~GeV), with a log-normal width $\sigma_L \sim 1.4{-}2.8$.
% These values are consistent with the luminosity function measured independently for MSPs in the Galactic Plane~\cite{Holst:2024fvb}, demonstrating that old stellar populations do not produce systematically fainter pulsars.
\blue{Building on an earlier study by Hooper and Linden~\cite{Hooper:2016rap}, in Amerio, Hooper, and Linden~\cite{Amerio:2024qor} we measured the gamma-ray luminosity function of MSPs across the Milky Way's globular cluster system, using the full Fermi-LAT dataset accumulated to date.}
\blue{The luminosity function we obtain is consistent with that measured independently for MSPs in the Galactic Plane~\cite{Holst:2024fvb}: we find no support for the hypothesis that old stellar populations produce systematically fainter pulsars.}

The implications for the GCE are direct.
If the excess is generated by MSPs with the luminosity function measured in globular clusters, Fermi-LAT's source catalogs should contain $N_\mathrm{MSP} \sim 17{-}37$ individually resolved pulsars from the Inner Galaxy population.
Only three MSP candidates have been identified in this region~\cite{Holst:2024fvb}, and systematic radio follow-up campaigns targeting unassociated Fermi sources have not so far uncovered a hidden population~\cite{Amerio:2024qor}.
% The GCE can therefore originate from MSPs only if they are significantly less luminous, on average, than those found in any other observed environment.
\blue{An MSP origin for the GCE would therefore require a pulsar population unlike any observed so far --- a requirement for which the hypothesis, as currently formulated, offers no explanation.}

This independent constraint connects directly to the recent findings of List et al.~\cite{List:2025qbx}.
As discussed in Section~\ref{sec:4.3.3}, their energy-dependent CNN analysis pushes the source-count distribution of any GCE point sources to extremely faint levels, requiring on the order of $2 \times 10^5$ sources --- each producing far less than one detected photon --- to account for the excess.
Such an ultra-faint population would naturally explain the absence of individually resolved pulsars.
% However, the globular cluster results remove the astrophysical motivation for its existence: if old age does not produce the required ultra-faint MSP populations, no known mechanism can generate $\mathcal{O}(10^5)$ sources of the necessary faintness in the Inner Galaxy.
\blue{However, the globular cluster results undermine the leading mechanism proposed to explain such a population: if old age does not produce ultra-faint MSPs, it is unclear what would generate $\mathcal{O}(10^5)$ sources of the necessary faintness in the Inner Galaxy.}
The tension identified from globular clusters and the tension from the energy-dependent SCD analysis thus converge from independent lines of evidence.
\blue{The remainder of this chapter is based on \cite{Amerio:2024qor} and presents the analysis in full.}

% \aure{An important comment. GLCs continue to form new MSPs through recycling, which would make younger and brighter MSPs. On the other hand, GLCs that inspiral into the GC get tidally disrupted, and cannot form new MSPs. So the MSP population in the GC is a mix of old MSPs from disrupted GLCs and younger MSPs formed in situ. This complicates the picture, but it does not change the main conclusion that the MSP population in the GC is old. }

% Hi Dan, hi Tim,

% Hope you are both doing well.

% I am currently writing my PhD thesis, and I'm working on the introduction to the problem of the GCE and the millisecond pulsar explanation. While reviewing the literature and going over our paper (on the MSP gamma-ray luminosity function in globular clusters), a question came up regarding the relative ages of these MSP populations, and I wanted to get your thoughts on how it might impact our conclusions.

% In our analysis, we studied surviving globular clusters to measure the gamma-ray luminosity function of their MSPs, motivated by the fact that both GLCs and the Galactic Bulge contain comparably old stellar populations. We argued that if the GCE is generated by pulsars with the same luminosity function we derived for the GLCs, then Fermi should have already detected roughly 17 to 37 bright MSPs close to the GC.

% However, while the stellar ages of these two environments are comparable, this doesn't guarantee that the ages of their respective MSP populations are the same. I am wondering if we would need to account for a divergence in the MSP formation history between surviving GLCs and those that deposited their pulsars into the GC.

% If I understand correctly, surviving GLCs are continually "recycling" new, bright MSPs through ongoing dynamical stellar encounters. But for the MSPs currently in the GC—which were likely deposited there via the dynamical friction and tidal disruption of inspiraling GLCs—this isn't the case. When those clusters were tidally disrupted by the Galaxy's deep gravitational well, they lost their extreme core density, leading to a sudden interruption in the dynamical formation of new MSPs.

% Doesn't this imply that the MSP population in the GC is potentially older (since their recycling was halted billions of years ago) compared to that of GLCs where new MSPs might still be formed today?

% If the GC MSPs have just been aging and losing rotational energy (dimming) without any new MSPs being formed to replace them, their gamma-ray luminosity function could still be systematically fainter than the one we measured in our paper.

% How heavily do you think this influences our hypothesis? Does this physical difference in the recycling history potentially introduce a loophole that could hide a fainter population from detection, meaning our estimate of 17-37 detectable MSPs might be too high?

% I'd love to hear your perspective on this, especially as I frame the discussion of our results in my thesis.

% Thank you very much,
% Best,

\aure{all the rest comes directly from the paper}